% Akathist booklet — XeLaTeX source scaffold.
% Fill the body from source.html (the formatting-preserving ground truth), then
% delete the guidance comments. Keep the preamble intact — it defines the Janson
% font and the formatting macros. See references/structure.md for how each source
% HTML element maps onto a macro.

\documentclass[11pt]{extarticle}
\usepackage[letterpaper,top=0.7in,bottom=0.7in,left=0.75in,right=0.75in]{geometry}
\usepackage{fontspec}
% Only the "55 Roman" OTF is installed, so synthesize bold + italic from it.
\setmainfont{Janson Text LT Std}[
  Numbers={OldStyle,Proportional},
  BoldFont={Janson Text LT Std},       BoldFeatures={FakeBold=2.2},
  ItalicFont={Janson Text LT Std},     ItalicFeatures={FakeSlant=0.18},
  BoldItalicFont={Janson Text LT Std}, BoldItalicFeatures={FakeBold=2.2,FakeSlant=0.18},
]
\usepackage{parskip}
\setlength{\parskip}{0.8\baselineskip plus 2pt}   % open spacing between paragraphs
\usepackage{ragged2e}       % \RaggedRight — even word spacing, no stretched gaps
\usepackage{needspace}      % keep a heading with the start of its stanza
\setlength{\emergencystretch}{2em}
% No hyphenation: keep every word whole so the exact-text QA gate never sees a
% word split across a line break.
\hyphenpenalty=10000
\exhyphenpenalty=10000
\pagestyle{plain}          % centered page number in the footer (aids booklet order)
\frenchspacing

% --- Formatting mirrors the source HTML exactly -----------------------------
% Apply weight/slant ONLY where the source marks it — via an element (h1/h2/h3
% are headings, <strong>/<b> bold, <em>/<i> italic) OR a CSS class (e.g.
% class="italic", a class containing "font-bold"). Body text, salutations, and
% refrains stay plain, because the source renders them as plain paragraph text.
% Never add bold or italic the source does not have.
\newcommand{\akatitle}[1]{\begin{center}\fontsize{24}{28}\selectfont\bfseries #1\end{center}}
% Optional dedication / commemoration line under the title, if the source has one.
\newcommand{\akasubtitle}[1]{\begin{center}\large #1\end{center}}
% Section heading, from an <h1>/<h2>/<h3>. Kept with the start of its stanza:
% \needspace forces a page break *before* the heading if too few lines remain
% (so a heading never dangles at the foot of a page); \nopagebreak glues it to
% the body paragraph that follows.
\newcommand{\heading}[1]{%
  \par\needspace{4\baselineskip}%
  \bigskip
  {\centering\bfseries #1\par}%
  \nopagebreak}
% Rubric — text the source sets in italics (e.g. class="italic"): a spoken or
% action instruction such as "(Repeat thrice)". Use ONLY when the source
% italicizes it; otherwise the line is a plain paragraph.
\newcommand{\rubric}[1]{{\itshape #1}\par}

\begin{document}
\fontsize{18}{24}\selectfont
\RaggedRight

% ---- TITLE (from the <h1>) ------------------------------------------------
\akatitle{% REPLACE with the <h1> text, e.g. Akathist to the Holy Martyr N.N.
}
\medskip

% ---- BODY -----------------------------------------------------------------
% One block per source element, in document order:
%   <h3>Kontakion 1</h3>      -> \heading{Kontakion 1}
%   <p>…</p>                  -> a plain paragraph (blank line before and after)
%   <p>a<br>b<br>c</p>        -> ONE paragraph, lines joined with \\  (the
%                                salutations and their closing refrain are one
%                                such paragraph — the refrain is just its last line)
%   <p class="italic">…</p>   -> \rubric{…}
%
% Worked example — replicate this shape for every Kontakion and Ikos:

\heading{Kontakion 1}
% REPLACE: the Kontakion 1 paragraph. Its refrain is inline, exactly as the source has it.

\heading{Ikos 1}
% REPLACE: the Ikos 1 lead paragraph.

% REPLACE: the salutations + closing refrain as ONE paragraph, one \\ per line.
% Everything here is plain — no bold, no hanging indent — matching the source:
Rejoice, % ...! \\
Rejoice, % ...! \\
Rejoice, O holy N.N., % ...  (closing refrain — plain, just the last line, no \\)

% ... continue through Kontakion 13, then any \rubric the source italicizes, then
% the repeated Ikos 1 and Kontakion 1 as the source directs.

\end{document}
